% =====================================================================
%  Chapter 3 — The Bullet Cluster and the Failure of Baryons
% =====================================================================

\chapter{The Bullet Cluster and the Failure of Baryons}
\label{ch:bullet}

Of all the evidence for dark matter, one object is more informative than the
rest. The merging cluster 1E~0657-56, known as the Bullet Cluster, was
analysed in 2006 with a conclusion that is hard to overstate: the dominant
mass component of a galaxy cluster is not the stuff that makes up stars and
gas, and it does not interact with that stuff when the cluster collides.

\section{What was observed}The Bullet is a textbook merging event: a smaller
subcluster, the "bullet", has passed through a larger subcluster at high
velocity, roughly 3{,}000~km\,s$^{-1}$ at the time of the observation, and
is now passing through the other side.

The three components were measured separately and they do not overlap.

\begin{itemize}
\item \textbf{The galaxies} --- the visible, ordinary stellar component ---
sit in one place, and are observed to be moving in a way consistent with two
separate pre-collision groups.
\item \textbf{The gas} --- the X-ray-emitting intracluster medium, which is
the dominant baryonic mass --- has been rammed aside and heated. The
collisional shock has pushed it away from the galaxies, and it is radiating
copiously in X-rays as it cools. The gas is observed to be in the process of
separating from the galaxies.
\item \textbf{The total mass} --- mapped by gravitational lensing of
background galaxies, and independently by the motions of the cluster's
galaxies --- sits at a third location, and it is nearly co-located with the
galaxies rather than with the gas.
\end{itemize}

The lensing mass is roughly 2.5 times the mass of the stars and the hot gas
combined. And the offset between the lensing mass and the gas is the
informative part: the gas is the only component that experienced a
collisional shock, and it is the only component that has moved away.

\section{What this rules out}Any candidate for the missing mass that is electromagnetically
interacting, or that has a substantial cross section for collision with
ordinary gas, is severely damaged by this observation. A baryonic solution
would need the hot gas of the cluster to be removed or made invisible in the
collision region, which is where the gas is observed to be, in the act of
cooling and radiating.

The conclusion Clowe and collaborators drew was correspondingly strong: the
dominant mass component is effectively collisionless and electromagnetically
inert, which is the definition of dark matter that the rest of this field
uses.

It is worth being precise about what the Bullet Cluster does \emph{not}
establish. It does not establish that dark matter is made of particles.
It does not exclude a light boson, a primordial black hole, a field, or a
population of compact objects. It constrains the interaction channels
extremely tightly and says nothing about the mass. A dark matter candidate
with a self-interaction cross section comparable to or larger than the
cluster's dark matter density times a velocity dispersion --- that is, of
order a few times $\mathrm{cm}^2/\mathrm{g}$ --- would produce a dark-matter
distribution offset from the galaxies, and is constrained by this cluster and
by others.

\section{The baryon budget problem}The Bullet Cluster is the sharpest version of a longer problem.
The baryonic catastrophe, as it was called, is the mismatch between the
density of baryons in the early universe predicted by nucleosynthesis and
the CMB, and the density of baryons actually observed in galaxies and
clusters today. On the largest scales the two agree: rich clusters are
observed to contain about 15 to 18 per cent of their mass in baryons, close
to the cosmic mean of about 16 per cent.

On smaller scales the agreement fails. Groups and low-mass galaxy clusters
are observed to be strongly baryon-depleted, and the standard account is
that their baryons have been pushed out into the surrounding medium by
feedback. The account is plausible and hard to verify, and it has a
consequence that matters for the rest of the book: if a large fraction of
baryons can be hidden in some environments, then a general argument that
``the baryons are not enough'' must be made carefully, environment by
environment. The Bullet Cluster is the environment where the argument is
airtight.

\section{What a modified theory of gravity has to reproduce}Modified gravity proposals -- most famously
MOND and its relativistic successors -- reproduce galactic rotation curves
extremely well, in some cases better than a dark matter halo fits, and this
is their strongest empirical result. They are not competitive with the
Bullet Cluster.

The requirement is straightforward. A modified force law that scales
differently from Newtonian gravity in low-acceleration environments must
still produce the observed lensing mass offset from the gas in a collision
where the baryons experience a shock and the dark component does not. That
requirement is met by no modified theory currently in the literature. It is
the single strongest reason the field treats dark matter as a real component
rather than as a placeholder for modified gravity.

Chapter~\ref{ch:modified} returns to the modified-gravity case in full and
is careful to distinguish what these theories do well from the ceilings they
have not escaped.

\begin{keytakeaways}
\item The Bullet Cluster separates three components of a merging cluster and
shows that the dominant mass component does not share the position of the
collided gas: it is effectively collisionless and electromagnetically inert.
\item This excludes the whole class of baryonic and electromagnetically
interacting explanations, and it constrains dark matter self-interactions to
be small.
\item The Bullet Cluster says nothing about whether dark matter is a
particle. It constrains interaction channels, not mass.
\item No modified theory of gravity currently reproduces the lensing mass
offset, which is the strongest reason the field treats dark matter as real.
\end{keytakeaways}
